\documentclass[11pt,a4paper]{article}

\usepackage[T1]{fontenc}
\usepackage[utf8]{inputenc}
\usepackage{lmodern}
\usepackage{amssymb,amsmath,amsthm,mathtools}
\usepackage[margin=2.5cm]{geometry}
\usepackage{parskip}
\usepackage[hidelinks]{hyperref}
\usepackage{microtype}

% Theorem environments (arXiv / journal style)
\theoremstyle{plain}
\newtheorem{theorem}{Theorem}[section]
\newtheorem{proposition}[theorem]{Proposition}
\newtheorem{lemma}[theorem]{Lemma}
\newtheorem{corollary}[theorem]{Corollary}

\theoremstyle{definition}
\newtheorem{definition}[theorem]{Definition}
\newtheorem{example}[theorem]{Example}
\newtheorem{remark}[theorem]{Remark}

% Notation
\newcommand{\Type}{\mathbf{Type}}
\newcommand{\id}{\mathrm{id}}
\newcommand{\Cont}{\mathbf{Cont}}
\newcommand{\Shape}{\mathrm{Shape}}
\newcommand{\Pos}{\mathrm{Pos}}
\newcommand{\Ext}{\mathrm{Ext}}
\newcommand{\seq}{\mathbin{\triangleleft}}

\title{Four Monoidal Structures on Containers:\\
  Coproduct, Product, and Extension Comparisons}
\author{Extracted from a machine-checked development}
\date{}

\begin{document}
\maketitle

\begin{abstract}
The category $\Cont$ of containers carries (at least) four monoidal
structures of interest: the coproduct~$+$, the categorical product~$\times$,
the Dirichlet tensor~$\otimes$, and the sequential operator~$\seq$. We
complete the account of the first two: bifunctoriality, structural
isomorphisms, naturality, and the pentagon and triangle coherence laws.
We further prove the comparison isomorphisms that make the extension
functor $\Ext(-)$ strong monoidal for the coproduct and the product,
\begin{align*}
  \Ext(C+D)\,X &\;\cong\; \Ext\,C\,X \;\oplus\; \Ext\,D\,X, \\
  \Ext(C \times D)\,X &\;\cong\; \Ext\,C\,X \;\times\; \Ext\,D\,X,
\end{align*}
together with the unit comparisons and the shared identity comparison
$\Ext\,y\,X \cong X$. The product/Dirichlet contrast is emphasised:
product and Dirichlet share shapes $S \times T$ and differ only in
positions ($P(s)\oplus Q(t)$ versus $P(s)\times Q(t)$), yet the extension
of the product is the pointwise product of functors by the exponential
law $X^{A+B} \cong X^A \times X^B$. The proofs of all stated results have
been machine-checked; we record only the mathematical content.
\end{abstract}

\tableofcontents
\bigskip

%======================================================================
\section{Introduction}
%======================================================================

Containers present polynomial endofunctors. The category $\Cont$ of
containers and their morphisms admits several tensors, each with a
different reading on extensions:
\begin{center}
\begin{tabular}{llll}
structure & tensor of $(S,P)$, $(T,Q)$ & unit & $\Ext$ pointwise? \\
\hline
coproduct $+$ & $(S \oplus T,\; \mathsf{elim}(P,Q))$ & $0$ & yes \\
product $\times$ & $(S \times T,\; (s,t) \mapsto P(s)\oplus Q(t))$ & $1$ & yes \\
Dirichlet $\otimes$ & $(S \times T,\; (s,t) \mapsto P(s)\times Q(t))$ & $y$ & no \\
sequential $\seq$ &
  $\bigl(\sum_s (P(s)\to T),\; \ldots\bigr)$ & $I{=}y$ & no
\end{tabular}
\end{center}
The sequential and Dirichlet structures are treated in companion notes.
Here we supply the monoidal data for~$+$ and~$\times$, and the extension
comparisons that make $\Ext(-)$ strong monoidal for both.

\paragraph{The product/Dirichlet contrast.}
Product and Dirichlet have the same shapes $S \times T$. They differ only
in positions: a sum versus a product. Yet the \emph{extensions} invert
the reading. The container whose positions are a sum has extension the
pointwise product of functors, because a labelling of $P(s)\oplus Q(t)$
\emph{is} a pair of labellings. The Dirichlet positions $P(s)\times Q(t)$
admit no such splitting, and $\Ext(C \otimes D)$ is not any pointwise
combination of $\Ext\,C$ and $\Ext\,D$.

\paragraph{What this note claims.}
\begin{enumerate}
  \item $(\Cont, +, 0)$ is monoidal.
  \item $(\Cont, \times, 1)$ is monoidal.
  \item $\Ext(C+D)\,X \cong \Ext\,C\,X \oplus \Ext\,D\,X$ and
        $\Ext\,0\,X \cong \emptyset$.
  \item $\Ext(C \times D)\,X \cong \Ext\,C\,X \times \Ext\,D\,X$ and
        $\Ext\,1\,X \cong 1$.
  \item $\Ext\,y\,X \cong X$ (shared unit comparison for $\otimes$
        and~$\seq$).
\end{enumerate}
None of the comparison maps is an identity: each is a genuine
isomorphism, which is exactly what ``strong, not strict'' means.

We work throughout in a single universe of types. Proofs are indicated by
the phrase ``by lean.''

%======================================================================
\section{Background and units}
%======================================================================

\begin{definition}[Container and extension]\label{def:container}
A \textbf{container} $C = (S,P)$ consists of a shape type~$S$ and a
position family~$P$. Its \textbf{extension} at a type~$X$ is
\begin{equation}
  \Ext\,C\,X \;=\; \sum_{s \,:\, S} \bigl(P(s) \to X\bigr).
\end{equation}
Morphisms of containers are shape maps with contravariant position maps;
we write $\Cont$ for the category.
\end{definition}

\begin{definition}[Empty and unit containers]\label{def:zero-one}
\begin{enumerate}
  \item The \textbf{empty container} $0$ has shape type $\emptyset$
        (and, vacuously, no positions). Its extension is the constantly
        empty functor. It is the unit for the coproduct and the initial
        object of~$\Cont$.
  \item The \textbf{one-shape, no-position container} $1$ has shape type
        $1$ (a singleton) and empty position set. Its extension is the
        constantly-one functor---the polynomial $y^0$. It is the unit for
        the categorical product and the terminal object of~$\Cont$.
\end{enumerate}
Note that $1 = y^0$ and $y = y^1$: the former is the unit for~$\times$,
the latter for~$\otimes$ and~$\seq$.
\end{definition}

We take as given the coproduct $C + D$ and product $C \times D$ of
containers (with their universal properties): shapes of the coproduct are
$S \oplus T$ with positions $\mathsf{elim}(P,Q)$; shapes of the product
are $S \times T$ with positions $P(s) \oplus Q(t)$.

%======================================================================
\section{The coproduct monoidal structure}
%======================================================================

\begin{definition}[Action on morphisms]\label{def:coprod2}
Given morphisms $\mu : C \to C'$ and $\nu : D \to D'$, the
\textbf{coproduct action}
\[
  \mu +_2 \nu \;:\; C + D \;\longrightarrow\; C' + D'
\]
is obtained from the coproduct's universal property: it is the copairing
of $\mu$ then the left injection and $\nu$ then the right injection.
\end{definition}

\begin{proposition}[Functor laws for $+$]\label{prop:coprod2-laws}
\begin{enumerate}
  \item $\id_C +_2 \id_D = \id_{C+D}$;
  \item $(\mu \mathbin{;}\mu') +_2 (\nu \mathbin{;}\nu')
        = (\mu +_2 \nu) \mathbin{;} (\mu' +_2 \nu')$.
\end{enumerate}
\end{proposition}

\begin{proof}
By lean.
\end{proof}

\begin{definition}[Structural isomorphisms for $+$]\label{def:coprod-struct}
\begin{enumerate}
  \item The \textbf{associator}
        $(C+D)+E \xrightarrow{\sim} C+(D+E)$ is $\oplus$-reassociation
        on shapes and the identity on positions in every branch (the
        positions of an injected shape \emph{are} the positions of the
        summand).
  \item The \textbf{left unitor} $0 + C \xrightarrow{\sim} C$ and
        \textbf{right unitor} $C + 0 \xrightarrow{\sim} C$ collapse the
        empty summand.
\end{enumerate}
\end{definition}

\begin{proposition}[Naturality for $+$]\label{prop:coprod-nat}
The associator and both unitors of the coproduct are natural in all
arguments.
\end{proposition}

\begin{proof}
By lean.
\end{proof}

\begin{theorem}[Coherence for $+$]\label{thm:coprod-coherence}
The pentagon and triangle identities hold for the coproduct tensor.
\end{theorem}

\begin{proof}
By lean.
\end{proof}

\begin{theorem}[$(\Cont, +, 0)$ is monoidal]\label{thm:coprod-monoidal}
The category of containers, equipped with the coproduct as tensor and the
empty container as unit, is a monoidal category.
\end{theorem}

\begin{proof}
By lean.
\end{proof}

\begin{remark}
In a formal development that admits only one canonical monoidal instance
on~$\Cont$, the sequential structure is preferred and the coproduct is
packaged as named data rather than a global instance. Mathematically
there is no conflict: four monoidal structures on one category is the
point.
\end{remark}

%======================================================================
\section{The product monoidal structure}
%======================================================================

\begin{definition}[Action on morphisms]\label{def:prod2}
Given morphisms $\mu : C \to C'$ and $\nu : D \to D'$, the
\textbf{product action}
\[
  \mu \times_2 \nu \;:\; C \times D \;\longrightarrow\; C' \times D'
\]
is obtained from the product's universal property: it is the pairing of
the first projection then~$\mu$ with the second projection then~$\nu$.
\end{definition}

\begin{proposition}[Functor laws for $\times$]\label{prop:prod2-laws}
\begin{enumerate}
  \item $\id_C \times_2 \id_D = \id_{C \times D}$;
  \item $(\mu \mathbin{;}\mu') \times_2 (\nu \mathbin{;}\nu')
        = (\mu \times_2 \nu) \mathbin{;} (\mu' \times_2 \nu')$.
\end{enumerate}
\end{proposition}

\begin{proof}
By lean.
\end{proof}

\begin{definition}[Structural isomorphisms for $\times$]\label{def:prod-struct}
\begin{enumerate}
  \item The \textbf{associator}
        $(C \times D) \times E \xrightarrow{\sim} C \times (D \times E)$
        is product reassociation on shapes and (backward) sum
        reassociation on positions: positions of the product are sums,
        so the associator must reassociate those sums in the opposite
        direction from the shapes.
  \item The \textbf{left unitor} $1 \times C \xrightarrow{\sim} C$ and
        \textbf{right unitor} $C \times 1 \xrightarrow{\sim} C$: the
        unit shape contributes no positions, so positions of the product
        are just those of~$C$ (injected into the appropriate summand).
\end{enumerate}
\end{definition}

\begin{proposition}[Naturality for $\times$]\label{prop:prod-nat}
The associator and both unitors of the product are natural in all
arguments.
\end{proposition}

\begin{proof}
By lean.
\end{proof}

\begin{theorem}[Coherence for $\times$]\label{thm:prod-coherence}
The pentagon and triangle identities hold for the product tensor.
\end{theorem}

\begin{proof}
By lean.
\end{proof}

\begin{theorem}[$(\Cont, \times, 1)$ is monoidal]\label{thm:prod-monoidal}
The category of containers, equipped with the categorical product as
tensor and the one-shape no-position container as unit, is a monoidal
category.
\end{theorem}

\begin{proof}
By lean.
\end{proof}

\begin{remark}
Shapes multiply, but positions add. That is the subtlety of the product
tensor, and the reason its extension comparison (Section~\ref{sec:ext})
is the exponential law rather than a mere unpacking of definitions.
\end{remark}

%======================================================================
\section{Extension as a monoidal functor}
\label{sec:ext}
%======================================================================

For a monoidal functor one needs, for each monoidal structure on the
source, a tensor on the target (endofunctors of $\Type$) and comparison
maps. Here:
\begin{itemize}
  \item $+$ maps to the pointwise coproduct of functors, $0$ to the
        constantly empty functor;
  \item $\times$ maps to the pointwise product of functors, $1$ to the
        constantly-one functor;
  \item $\seq$ maps to composition of functors (companion note);
  \item $\otimes$ admits \emph{no} pointwise tensor on the target
        (companion note on the Dirichlet structure).
\end{itemize}
Each comparison below is a genuine isomorphism of types.

\subsection{Coproduct comparisons}

\begin{definition}[Coproduct extension comparison]\label{def:coprod-ext}
The forward and backward maps
\begin{align}
  \Ext(C+D)\,X
  &\;\longrightarrow\;
  \Ext\,C\,X \;\oplus\; \Ext\,D\,X, \\
  \Ext\,C\,X \;\oplus\; \Ext\,D\,X
  &\;\longrightarrow\;
  \Ext(C+D)\,X
\end{align}
send a structure whose shape is a left (resp.\ right) injection to the
corresponding summand, and conversely re-inject. The unit comparison
$\Ext\,0\,X \to \emptyset$ is elimination of the empty shape.
\end{definition}

\begin{theorem}[Coproduct comparisons are inverse]\label{thm:coprod-ext}
The maps of Definition~\ref{def:coprod-ext} are mutually inverse. In
particular
\begin{equation}
  \Ext(C+D)\,X \;\cong\; \Ext\,C\,X \;\oplus\; \Ext\,D\,X,
  \qquad
  \Ext\,0\,X \;\cong\; \emptyset.
\end{equation}
\end{theorem}

\begin{proof}
By lean.
\end{proof}

\subsection{Product comparisons}

\begin{definition}[Product extension comparison]\label{def:prod-ext}
The forward map
\[
  \Ext(C \times D)\,X
  \;\longrightarrow\;
  \Ext\,C\,X \;\times\; \Ext\,D\,X
\]
splits a labelling of $P(s) \oplus Q(t)$ along the two injections. The
backward map pairs the shapes and combines the labellings by case
analysis on the sum of positions.
\end{definition}

\begin{theorem}[Product comparisons are inverse]\label{thm:prod-ext}
The maps of Definition~\ref{def:prod-ext} are mutually inverse. In
particular
\begin{equation}
  \Ext(C \times D)\,X \;\cong\; \Ext\,C\,X \;\times\; \Ext\,D\,X,
  \qquad
  \Ext\,1\,X \;\cong\; 1.
\end{equation}
\end{theorem}

\begin{proof}
By lean.
\end{proof}

\begin{remark}[The exponential law]
This is the pleasing cross-check. The product container's positions are a
\emph{sum}, yet its extension is the pointwise \emph{product} of the
extensions. The reason is
\begin{equation}
  X^{A+B} \;\cong\; X^A \times X^B:
\end{equation}
a labelling $P(s) \oplus Q(t) \to X$ \emph{is} a pair of labellings. Currying
the sum turns the position-level $\oplus$ into a functor-level $\times$.
\end{remark}

\subsection{The shared unit $y$}

\begin{definition}[Identity container as representable]\label{def:y}
Write $y$ for the container with one shape and one position (also called
$I$ when viewed as the unit of~$\seq$). Its extension is the identity
functor.
\end{definition}

\begin{theorem}[Unit comparison for $y$]\label{thm:y-ext}
The maps
\begin{equation}
  \Ext\,y\,X \;\longrightarrow\; X,
  \qquad
  X \;\longrightarrow\; \Ext\,y\,X
\end{equation}
---read the label at the unique position, and place a value at the unique
position---are mutually inverse. Thus $\Ext\,y\,X \cong X$. This is the
unit comparison for both $\otimes$ and~$\seq$.
\end{theorem}

\begin{proof}
By lean.
\end{proof}

%======================================================================
\section{Discussion}
%======================================================================

\paragraph{What has been established.}
The coproduct and product are monoidal structures on~$\Cont$. The
extension functor is strong monoidal for both, via genuine (not identity)
comparison isomorphisms. The product/Dirichlet contrast is explained by
the exponential law: sum of positions yields product of extensions;
product of positions does not split.

\paragraph{Taxonomy of the four structures.}
\begin{center}
\begin{tabular}{llll}
structure & Day convolution of & unit & $\Ext$ pointwise? \\
\hline
product $\times$ & $(\mathbf{Set}, +, 0)$ & $1 = y^0$ & yes \\
Dirichlet $\otimes$ & $(\mathbf{Set}, \times, 1)$ & $y = y^1$ & no \\
coproduct $+$ & not a Day convolution & $0$ & yes \\
sequential $\seq$ & not a Day convolution & $I = y$ & no
\end{tabular}
\end{center}
(The Day-convolution reading of $\times$ and $\otimes$ follows Niu--Spivak;
the coproduct fails to preserve the initial object in each variable, and
$\seq$ is not cocontinuous in the right variable.)

\paragraph{Machine checking.}
Every definition and theorem statement above corresponds to a declaration
in a machine-checked formalisation. The phrase ``by lean'' indicates that
the proof has been verified by a proof assistant.

%======================================================================
\section*{Acknowledgements}
%======================================================================

The mathematical content follows the theory of containers of Abbott,
Altenkirch, and Ghani, and the monoidal taxonomy of polynomial functors.
The presentation here is a prose rendering of a formal development
intended for human review of the theorem statements.

\begin{thebibliography}{9}

\bibitem{abbott2005}
M.~Abbott, T.~Altenkirch, and N.~Ghani.
\newblock Containers: constructing strictly positive types.
\newblock \emph{Theoretical Computer Science}, 342(1):3--27, 2005.

\bibitem{abbott2003}
M.~Abbott, T.~Altenkirch, and N.~Ghani.
\newblock Categories of containers.
\newblock In \emph{FoSSaCS}, volume~2620 of \emph{LNCS}, pages 23--38.
\newblock Springer, 2003.

\bibitem{niu-spivak2023}
N.~Niu and D.~I.~Spivak.
\newblock Polynomial functors: a mathematical theory of interaction.
\newblock arXiv:2312.00990, 2023.

\bibitem{maclane1971}
S.~Mac~Lane.
\newblock \emph{Categories for the Working Mathematician}.
\newblock Springer, 1971.

\end{thebibliography}

\end{document}
